Să se proiecteze un automat secvențial sincron care controlează funcționarea unui aparat portabil de emisie-recepție (Walkie-Talkie).

Scopul automatului este de a facilita comunicarea bidirecțională între două dispozitive, permițând utilizatorului să transmită și să primească un flux audio.

Sistemul pornește într-o stare de așteptare (Idle/Listening), monitorizând semnalele de intrare. Comportamentul dispozitivului este dictat prin următoarele stări:

\begin{itemize}
    \item \textbf{Modul Normal (Receiving/Transmitting):} În starea implicită, dispozitivul primește un flux audio de la canalul său curent.
    Dacă se apasă butonul \texttt{PTT} (Push-to-Talk), automatul trece în modul de transmisie, activând microfonul și LED-ul de transmisie.
    La eliberarea butonului, revine în modul de recepție.
    
    \item \textbf{Modul Pairing:} La apăsarea butonului \texttt{PAIR}, dispozitivul intră într-o secvență de scanare, așteptând un răspuns (handshake), urmată de o stare de succes sau eroare.
    
    \item \textbf{Selectarea Canalului:} Dispozitivul permite comutarea între două canale salvate. Canalul activ este schimbat prin apăsarea \texttt{CH\_BTN}.
\end{itemize}
